% =============================================================================
% 5.3 全文总结
% 草稿版本：v0.1 (2026-05-05)
% 字数：约 1000 字
% 风格规则：v0.2 风格（无 \textbf / \textit / 引例括号 / 双引号 / 破折号）
% 引用：暂无新增引用
% 图表：
%   - 表 \reftab{tab:evidence-summary}    5 项实证一表概览
% 依赖宏包：booktabs / tabularx（主稿已加载）
% =============================================================================

\subsection{全文总结}
\label{sec:summary}

本文围绕基于 AI 智能体的气象预报数据检索与分析这一研究课题，针对开题
报告中提出的气象任务指令解析（Q1）、气象数据异构特征（Q2）、大语言
模型在数值计算与逻辑推理中的幻觉（Q3）三个研究问题，设计并实现了一套
完整的单智能体气象 Agent 系统，并以 5 类评测共 214 条用例进行了系统化
实证。完成的工作可归纳为五项贡献。

第一项贡献是基于 ReAct 范式的气象 Agent 单智能体架构。本文采用单
智能体加显式工具学习而非多智能体协作的工程取舍，把场景编排、出处
强制、出行场景多地点对齐等约束下沉到 System Prompt，避免了多智能体
之间的状态同步与额外延迟。整套架构在
\texttt{src/\bk{}agent/\bk{}react\_\bk{}agent.\bk{}py} 中以 LangGraph 与 LangChain 0.3
的 \texttt{create\_\bk{}agent} 实现，集成 9 个气象 API 工具加 2 个 RAG
工具，支持流式中间过程输出与多轮对话状态保持。

第二项贡献是工具学习驱动的异构数据检索机制。本文为和风天气与
Open-Meteo 两类异构 API 设计了统一的 Function Calling Schema，并通过
意图识别与参数补全双阶段管线把自然语言查询映射为结构化
\texttt{WeatherIntent} 对象。35 条评测的端到端通过率达到 91.4\%，
出行场景多地点角色识别准确率达到 100\%。分层架构本身即是容错机制，
单点大语言模型抖动可以由后置补全器在业务层吸收。

第三项贡献是双通路 RAG 架构。通路 A 混合检索把向量检索与 BM25 词法
检索以 0.9 与 0.1 加权融合，在 60 条评测集上 Top-1 命中率达到 85.0\%，
MRR 达到 0.908。通路 B 语义桥接由 5 类确定性分级器加 \texttt{grade\_\bk{}id}
硬链接构成，在 71 条评测集上 \texttt{grade\_\bk{}id} 准确率达到 100\%，
权威标准 \texttt{citation} 真实性达到 100\%，与大语言模型自由发挥的
基线在 \texttt{must\_\bk{}cite} 严格命中率上形成 74.6 个百分点的差距。
两条通路共同构成无重叠盲区的检索闭环。

第四项贡献是基于代码生成与受限沙箱的确定性统计分析机制。本文设计了
LLM 生成 \texttt{compute(data)} 顶层函数加白名单 builtins 加白名单
模块加墙钟超时加结构化错误捕获的轻量沙箱方案。在 18 条覆盖 6 类统计
模式的评测集上，代码执行档数值准确率达到 88.9\%，比大语言模型自由
直算高出 22.2 个百分点。失败模式分析进一步发现极值任务上自由直算
全军覆没（0\% 对比 66.7\%），计数任务上差距达 50 个百分点，揭示了
大语言模型在伴随状态维护与离散计数任务上的内在幻觉边界。

第五项贡献是覆盖模块层与集成层的多层评测体系。本文构建了 5 份独立
评测集共 214 条用例，覆盖意图识别、通路 A 检索、通路 B 桥接、代码
执行 4 个模块层评测加 1 个 30 条端到端集成层评测。集成层评测首次
以最终用户体验视角验证整套系统的端到端价值，揭示了集成衰减效应——
通路 B 单测 \texttt{citation} 真实性 100\% 在端到端集成后衰减到
26.7\%。5 组实证形成完整的证据链，详见 \reftab{tab:evidence-summary}。

\begin{table}[!htbp]
    \centering
    \caption{5 项评测的核心实证一表概览}
    \label{tab:evidence-summary}
    \song\fontsize{9pt}{12pt}\selectfont
    \renewcommand{\arraystretch}{0.95}
    \begin{tabular}{llrl}
        \toprule
        实证 & 评测对象 & 用例数 & 关键数据 \\
        \midrule
        意图理解 & 自然语言到 \texttt{WeatherIntent} & 35 & 端到端 91.4\% / 角色 100\% \\
        通路 A 检索 & 混合检索器 3 档 + 8 权重点 & 60 & Top-1 85.0\% / MRR 0.908 \\
        通路 B 桥接 & 4 档处理方式对照 & 71 & \texttt{grade\_\bk{}id} 100\% / \texttt{citation} 100\% \\
        代码执行 & 3 档消融对照 & 18 & 数值准确率 +22.2 pp \\
        端到端集成 & \texttt{full} 与 \texttt{ablated} 双档 & 30 & 通过率 +13.4 pp / 引用率 +26.7 pp \\
        \bottomrule
    \end{tabular}
\end{table}

5 组实证在数据规模、研究问题与实证视角上互不重叠，共同构成本文论点
的完整闭环。Q1 意图理解的 91.4\% 端到端通过率回答了用户提问能否被
正确解析；Q2 检索 Top-1 85\% 加桥接 \texttt{grade\_\bk{}id} 100\% 回答了
异构数据能否被可靠取出与解读；Q3 代码执行 +22.2 pp 加桥接
\texttt{citation} 真实性 100\% 回答了数值与引用层面的幻觉能否被有效
缓解；端到端 +13.4 pp 与 +26.7 pp 共同回答了集成层是否产生工程可用
的最终输出。本文据此完成了开题报告 §3.4 承诺的 5 项评测指标的全部
实证支撑。

